% !PRINTED 60
\weberpage{60}
%<<CONT

\[
 \alpha_1+\alpha_2+\cdots+\alpha_m=\nu
\]
ist, die \gesperrt{Derivirten $\nu$ter Ordnung der Function $\Phi$}.

Man stellt sie auch nach der in der Differentialrechnung gebräuchlichen Bezeichnungsweise so dar:
\begin{equation}
 D_{\alpha_1,\alpha_2,\ldots,\alpha_m}\Phi
 =
 {\partial^\nu\Phi\over
 \partial x_1^{\alpha_1}\partial x_2^{\alpha_2}\cdots\partial x_m^{\alpha_m}}.
\tag{5}
\end{equation}

Das Bildungsgesetz der Derivirten lässt sich in folgende Sätze zusammenfassen, wobei wir der Kürze wegen die Indices bei dem Zeichen $D$ weglassen.

\gesperrt{I. Ist $C$ eine Constante, so ist}
\[
 D(C\Phi)=CD\Phi.
\]
\gesperrt{II. Sind $\Phi$ und $\Psi$ irgend zwei Functionen, so ist}
\[
 D(\Phi+\Psi)=D\Phi+D\Psi.
\]
Beides folgt unmittelbar aus (4).

Wir können also leicht die derivirten Functionen allgemein bilden, wenn wir sie für den speciellen Fall kennen, in dem $\Phi$ ein Product von Potenzen ist, also wenn wir
\[
 D_{\alpha_1,\alpha_2,\ldots,\alpha_m}(x_1^{\mu_1}x_2^{\mu_2}\cdots x_m^{\mu_m})
\]
kennen, worin die $\mu$ beliebige, nicht negative Exponenten sind.

Nun ist aber
\[
 (x_1+\xi_1)^{\mu_1}
 =\sum^{\alpha_1}{\Pi(\mu_1)\over\Pi(\alpha_1)\Pi(\mu_1-\alpha_1)}
 \xi_1^{\alpha_1}x_1^{\mu_1-\alpha_1}
\]
und folglich:
\begin{equation}
\begin{aligned}
&(x_1+\xi_1)^{\mu_1}\cdots(x_m+\xi_m)^{\mu_m}\\
&=\sum^{\alpha_1\ldots\alpha_m}
 {\Pi(\mu_1)\cdots\Pi(\mu_m)x_1^{\mu_1-\alpha_1}\cdots x_m^{\mu_m-\alpha_m}
 \over
 \Pi(\mu_1-\alpha_1)\cdots\Pi(\mu_m-\alpha_m)\Pi(\alpha_1)\cdots\Pi(\alpha_m)}
 \xi_1^{\alpha_1}\cdots\xi_m^{\alpha_m}.
\end{aligned}
\tag{6}
\end{equation}
und es ergiebt also die Vergleichung mit (4) und (6)
\begin{equation}
\begin{aligned}
&D_{\alpha_1,\ldots,\alpha_m}(x_1^{\mu_1}\cdots x_m^{\mu_m})\\
&\quad =
{\Pi(\mu_1)\cdots\Pi(\mu_m)\over
 \Pi(\mu_1-\alpha_1)\cdots\Pi(\mu_m-\alpha_m)}
 x_1^{\mu_1-\alpha_1}\cdots x_m^{\mu_m-\alpha_m},
\end{aligned}
\tag{7}
\end{equation}
so lange $\alpha_1\leq\mu_1,\ldots,\alpha_m\leq\mu_m$. Dagegen ist
\begin{equation}
 D_{\alpha_1,\ldots,\alpha_m}(x_1^{\mu_1}\cdots x_m^{\mu_m})=0,
\tag{8}
\end{equation}
sobald einer der Indices $\alpha$ grösser ist als der entsprechende Exponent $\mu$.
